\section{Local Feature Extraction}
In this section we will explore a few methods for extracting local features from images.
In particular, we have already previously seen the \textbf{Good Features to Track} algorithm, in which we can extract
a set of corner points from an image, which are then tracked using the \textbf{Lucas-Kanade} method.
We will be exploring a few other methods for extracting local features, such as the \textbf{Histogram of Oriented Gradients (HOG)} and the \textbf{Scale-Invariant Feature Transform (SIFT)}.

\subsection{Histogram of Oriented Gradients (HOG)}
\subsubsection{Gradients}
We have already seen that a gradient is a vector that points in the direction of the greatest change in intensity. In image processing, 
the gradient is used to detect edges in an image, which are areas of high intensity change. Intensity and direction of edges are one of the
most important features usefull to carachterize an image.


\subsection{Scale-Invariant Feature Transform (SIFT)}
Had a lot of success in the past, since it is efficienct and easy to implement.